\completa{21}{Sergio}{5a.2: año, sexo y filtro agregado}
\senalar{Leer las cuatro cajas. Distinguir sexo F del umbral sobre personas.}
\decir{Este segundo ejemplo de cinco a nace de AÑO, SEXO y CANTIDAD. Preguntamos: ¿en qué años el total reportado de mujeres supera cinco mil personas? Queremos comprobar el umbral en cada año, manteniendo separadas las categorías de sexo antes de seleccionar F.}
\decir{El primer cuadro dice «Todos los registros». La flecha llega a «groupBy año, sexo»: cada pareja año y categoría de sexo forma un grupo. La siguiente conduce a «SUM de cantidad; sexo F y total mayor que cinco mil». Calculamos el acumulado y conservamos los grupos F cuyo total supera el umbral. La última flecha llega a «seis años; orden ascendente». Ascendente significa del año menor al mayor, de 2019 a 2024.}
\decir{En el código, agg reúne sum y count: personas es la suma de CANTIDAD y registros es el número de filas. Filter combina dos condiciones mediante el operador y: sexo igual a F y personas mayor que cinco mil. La igualdad de columnas en esta API se escribe con tres signos iguales. OrderBy anio aplica el orden cronológico.}
\decir{La salida muestra 2019 con dieciséis mil cuatrocientas cincuenta y una mujeres y once mil treinta y siete registros; 2020 con once mil ochocientas veintiocho y ocho mil setecientos veintitrés; 2024 con once mil doscientas setenta y cinco y ocho mil trescientos noventa y ocho. Los seis años cumplen el umbral, aunque la tabla visible muestra tres. El criterio se aplica al grupo, no a una fila aislada. Cambiaremos ahora las dos claves para observar región y periodo.}

\completa{22}{Sergio}{5a.3: dos variables derivadas y periodos desiguales}
\senalar{Señalar que región y periodo se crean en el esquema común. Leer las cuatro cajas y las cuatro filas visibles.}
\decir{La tercera agrupación de cinco a pregunta: ¿qué combinaciones de región y periodo superan diez mil personas reportadas? Nace de las columnas derivadas región y periodo, junto con la medida CANTIDAD. Buscamos acumulados grandes según esa clasificación territorial y temporal.}
\decir{El primer cuadro dice «Región y periodo: derivados del código». La flecha conduce a «groupBy: dos columnas». Formamos una combinación por región y periodo. La siguiente llega a «SUM de cantidad; total mayor que diez mil». Sumamos dentro de cada combinación y luego filtramos el acumulado. La última flecha termina en «cinco combinaciones; orden descendente». El código usa agg y as para producir personas, filter para el umbral y orderBy con desc para ubicar los mayores al inicio.}
\decir{En la salida, Costa de 2022 a 2024 suma treinta y siete mil doscientas cuatro personas; Costa de 2020 a 2021, veinticuatro mil seiscientas tres; Costa de 2019, diecisiete mil ochocientas cincuenta y nueve. Sierra de 2022 a 2024 suma quince mil novecientas ochenta y una. Son cuatro filas mostradas de cinco resultados. El mayor acumulado es Costa en el periodo final. Sin embargo, los periodos duran tres, dos y un año: su comparación no equivale a comparar promedios anuales. La siguiente consulta utilizará categorías originales de nacionalidad y edad.}

\completa{23}{Sergio}{5a.4: delimitar extranjeros antes de agrupar}
\senalar{Leer las cuatro cajas y contrastar el filtro inicial con el umbral posterior.}
\decir{Del cruce entre NACIONALIDAD y RANGO EDAD surge esta pregunta: ¿qué combinaciones de nacionalidad extranjera y rango de edad reúnen al menos treinta personas? Es el cuarto ejemplo de cinco a. Queremos identificar los grupos más numerosos del subconjunto extranjero declarado.}
\decir{El primer cuadro dice «Excluir PERÚ y NO ESPECIFICA». Este filtro actúa sobre las filas de entrada. La flecha lleva a «groupBy: nacionalidad y edad»: agrupamos únicamente las filas conservadas. La siguiente llega a «SUM de cantidad; total mayor o igual a treinta». Ahora el filtro actúa sobre grupos ya sumados. La última flecha conduce a «nueve combinaciones; orden descendente». Primero se ubican las de mayor número de personas.}
\decir{En el código, col referencia una columna; el operador de distinto descarta las dos etiquetas indicadas y el operador y exige ambas condiciones. Agg calcula sum, as nombra personas y filter aplica el mínimo de treinta. La salida muestra Venezuela de doce a diecisiete años con mil nueve personas; Venezuela de dieciocho a veintinueve con quinientas noventa y cinco; y Colombia de dieciocho a veintinueve con cincuenta y dos. El mayor volumen corresponde al primer grupo. No es una tasa respecto de cuántas personas de cada nacionalidad viven en el país. Después del volumen revisaremos una pregunta sobre promedios.}

\completa{24}{Sergio}{5b.1: promedio por registro y departamento}
\senalar{Seguir las cuatro cajas y explicar la diferencia entre promedio y total de la tabla.}
\decir{El punto cinco b también solicita un promedio. A partir de DEPARTAMENTO y CANTIDAD preguntamos: ¿cuál es el promedio de personas contabilizadas por registro en cada departamento? Esta pregunta es diferente de cuál departamento suma más personas. El promedio divide la suma de cantidades entre el número de filas del departamento.}
\decir{El primer cuadro dice «Registros por departamento». Su flecha llega a «groupBy: departamento». La siguiente conduce a «avg de cantidad; redondeo a tres decimales». Avg calcula la media y round fija las posiciones decimales de presentación. La última flecha lleva a «veintiséis departamentos; promedio descendente». Ordenamos por esa media, de mayor a menor. El código añade count para registros y sum para personas; as asigna nombres a las columnas calculadas.}
\decir{La salida muestra Callao con promedio uno coma cinco cero cinco y cuatro mil quinientas sesenta y cinco personas; Lima Metropolitana con uno coma cuatro ocho cinco y treinta y nueve mil cuatrocientas noventa; y Lambayeque con uno coma tres uno nueve y seis mil doscientas treinta y una. Callao encabeza la media, aunque Lima Metropolitana tiene mayor acumulado. La media describe la cantidad por cruce agregado, no cuántas personas participan de un mismo incidente. Veamos ahora cómo se detecta una falta de correspondencia entre tablas.}

\completa{25}{Sergio}{5c.2: qué conserva LEFT y qué significa un nulo}
\senalar{Identificar el catálogo como izquierda y los acumulados como derecha antes de seguir las flechas.}
\decir{El segundo JOIN de cinco c pregunta: ¿qué ubigeos del catálogo construido con el dataset no tienen registros asociados? Queremos mostrar cómo conservar una lista de referencia y detectar sus elementos sin coincidencia. La referencia es dimGeografia, a la izquierda, y los acumulados por ubigeo quedan a la derecha.}
\decir{El primer cuadro dice «Catálogo derivado del mismo dataset». La flecha llega a «LEFT OUTER JOIN: conservar todos sus ubigeos». Left conserva cada clave de la tabla izquierda, incluso si no encuentra coincidencia en la derecha; outer indica esa conservación de filas sin pareja. La siguiente flecha conduce a «personas punto isNull: buscar faltantes». Un null representa ausencia de valor; aquí aparecería cuando un ubigeo del catálogo no encuentra su acumulado. IsNull comprueba esa ausencia. La última flecha termina en «cero sin registros; mil quinientos seis con registros».}
\decir{En el fragmento se ve join con la clave ubigeo y el tipo left, seguido del filtro de personas nulas. La salida repite los tres controles: mil quinientos seis ubigeos de catálogo, todos con registros, y cero sin registros. Este cero verifica correspondencia interna. Como el catálogo procede de las propias denuncias, no prueba que estén todos los distritos del país. Para detectar ausencias respecto de un universo independiente necesitaremos una referencia construida fuera de las filas observadas, como el calendario de la siguiente consulta.}

\completa{26}{Sergio}{5c.3: RIGHT con un calendario independiente}
\senalar{Señalar el calendario a la derecha y distinguir coalesce como función de valores.}
\decir{La tercera unión pregunta: ¿hay meses sin registros entre enero de 2019 y diciembre de 2024? Aquí sí podemos construir la referencia completa sin depender de las denuncias: seis años por doce meses son setenta y dos combinaciones. Los acumulados están a la izquierda y ese calendario a la derecha.}
\decir{El primer cuadro dice «groupBy año, mes». Sumamos las cantidades de cada combinación observada. La flecha conduce a «RIGHT OUTER JOIN: conservar calendario». Right conserva todas las filas de la derecha, aunque un mes no tenga coincidencia en los acumulados. La siguiente llega a «coalesce a cero; filtrar total igual a cero». En esta expresión, coalesce elige el primer valor no nulo: mantiene el acumulado si existe y utiliza cero si falta. Después buscamos esos ceros. La última flecha termina en «setenta y dos meses presentes; cero meses sin registros».}
\decir{El código utiliza join de tipo right con las claves año y mes. WithColumn crea o sustituye la columna ajustada, y lit cero aporta el valor constante. Esta función coalesce sobre columnas no es el método coalesce uno de la exportación, que modifica particiones. La tabla confirma setenta y dos meses conservados, ninguno sin registros y una suma de ciento veintiún mil cuatrocientas veintisiete personas. La consulta controla presencia mensual; no garantiza que cada denuncia posible se haya registrado. Pasaremos ahora a conservar claves de ambos lados de una unión.}

\completa{27}{Sergio}{5c.4: FULL OUTER y el universo de edades conocidas}
\senalar{Leer las cuatro cajas; señalar por qué el denominador difiere del de 3c.1.}
\decir{El cuarto JOIN pregunta: ¿qué departamentos aparecen solo entre menores o solo entre adultos con edad conocida? RANGO EDAD permite separar dos subconjuntos; DEPARTAMENTO permite compararlos. Excluimos NO ESPECIFICA, por lo que esta consulta tiene un universo diferente del porcentaje de tres c uno.}
\decir{El primer cuadro dice «Filtrar edades conocidas; dos agregados». Sumamos menores por departamento y, por separado, adultos por departamento. La flecha conduce a «FULL OUTER JOIN: clave departamento». Full outer conserva todas las claves de ambos lados: si un departamento aparece solo en uno, no desaparece del resultado. La siguiente flecha llega a «coalesce a cero; sumar y calcular porcentaje». Reemplazamos cantidades ausentes por cero, sumamos menores y adultos y calculamos la participación de menores sobre ese total conocido. La última flecha conduce a «veintiséis departamentos; ambos grupos presentes».}
\decir{En el código, full guion bajo outer selecciona ese comportamiento; withColumn crea los totales y round redondea el porcentaje. La salida muestra Lima Metropolitana con diecinueve mil setecientos cincuenta y un menores y diecinueve mil doscientos ochenta y un adultos; Junín con cuatro mil quinientos veinticinco y dos mil seiscientos ochenta y ocho; Cusco con cuatro mil doscientos noventa y nueve y dos mil setecientos sesenta y dos.}
\decir{Los veintiséis departamentos tienen ambos grupos. Para Lima, menores dividido entre menores más adultos da cincuenta coma sesenta por ciento. En tres c uno, el mismo numerador se divide entre el total que incluye edad no especificada y da cincuenta coma cero dos. Cambiar el denominador cambia la pregunta; esa diferencia no es una contradicción. La siguiente ficha mostrará cómo crear categorías mediante condiciones.}

\completa{28}{Sergio}{5d.1: interpretar las etiquetas de when y otherwise}
\senalar{Seguir las cuatro cajas. Leer las etiquetas del código y explicar sus límites.}
\decir{Esta consulta nace de la variación de CANTIDAD y ejemplifica las funciones del punto cinco d. Preguntamos: ¿cuántas filas y personas corresponden a cada tramo de CANTIDAD definido por el código? Los tramos sirven para describir la distribución de conteos. No son categorías oficiales de gravedad.}
\decir{El primer cuadro dice «cantidad por fila». La flecha llega a «when y otherwise: uno a dos, tres a cuatro, cinco o más». When significa cuando: si cantidad es al menos cinco se asigna FOCO; si no, pero es al menos tres, se asigna ALTO; otherwise, en los demás casos de estos datos positivos, se asigna INDIVIDUAL. Importa comprobar primero el umbral más alto. WithColumn crea la columna llamada severidad con esa etiqueta.}
\decir{La siguiente flecha conduce a «groupBy severidad; count, sum y avg». Agrupamos por la nueva etiqueta, contamos filas, sumamos personas y calculamos el promedio por fila. La última flecha llega a «tres categorías; noventa y ocho mil doscientas treinta y seis, catorce mil tres y nueve mil ciento ochenta y ocho». Son sus sumas de personas, no sus conteos de filas.}
\decir{La tabla detalla INDIVIDUAL, uno a dos, con ochenta y ocho mil ciento ochenta y dos filas y noventa y ocho mil doscientas treinta y seis personas; ALTO, tres a cuatro, con cuatro mil doscientas cuarenta y nueve filas y catorce mil tres personas; FOCO, cinco o más, con mil cuatrocientas filas y nueve mil ciento ochenta y ocho personas. Aunque la columna se llame severidad y una etiqueta diga INDIVIDUAL, los datos no permiten afirmar gravedad ni incidentes individuales o grupales. Los nombres son decisiones del código sobre conteos agregados. Ahora calcularemos participaciones por periodo.}

\completa{29}{Sergio}{5d.2: numerador por periodo y denominador constante}
\senalar{Leer las cuatro cajas; señalar que lit introduce un valor fijo y round solo redondea.}
\decir{De la columna periodo nace la pregunta: ¿qué participación tiene cada periodo sobre el total de personas reportadas? Es otra aplicación de funciones del punto cinco d. Necesitamos un numerador que cambia por periodo y un denominador nacional común.}
\decir{El primer cuadro dice «Obtener total nacional: agg y sum». La agregación de todo el DataFrame obtiene ciento veintiún mil cuatrocientas veintisiete. La flecha conduce a «groupBy periodo: suma por periodo». Cada periodo recibe su propio acumulado. La siguiente llega a «lit total nacional; round de cien por parte dividido entre total». Lit convierte el número fijo en una expresión de columna utilizable en la fórmula. Round redondea el porcentaje; no cambia su denominador. La última flecha llega a «tres participaciones; suman cien por ciento».}
\decir{En el código, col personas referencia el numerador y withColumn añade pct nacional. Se añade también personas por registro dividiendo la suma de personas entre count de registros, con tres decimales. La salida de 2019 es veintiséis mil trescientas treinta y una personas, veintiuno coma sesenta y ocho por ciento; 2020 a 2021 suma treinta y ocho mil seiscientas veintidós, treinta y uno coma ochenta y uno por ciento; 2022 a 2024 suma cincuenta y seis mil cuatrocientas setenta y cuatro, cuarenta y seis coma cincuenta y uno por ciento.}
\decir{El periodo final pesa más dentro del acumulado, pero dura tres años frente a dos y uno de los demás. Esa participación no demuestra mayor intensidad anual ni un efecto causal de la pandemia. Si quisiéramos intensidad comparable, tendríamos que definir otra medida que tenga en cuenta la duración. La siguiente consulta construye una unidad temporal más detallada: el trimestre de cada año.}

\completa{30}{Sergio}{5d.4: construir una fecha sin inventar un día observado}
\senalar{Seguir las cuatro cajas y relacionar concat, lpad y to date con la fecha convencional.}
\decir{AÑO y MES permiten preguntar: ¿cuántas personas se reportan por trimestre al construir una fecha con año y mes? Esta consulta del punto cinco d muestra funciones temporales. Necesitamos una fecha utilizable por Spark, aunque el archivo no informe el día. Por eso fijamos convencionalmente el día uno.}
\decir{El primer cuadro dice «Año y mes originales». La flecha lleva a «concat ws más lpad; to date». Concat ws une textos con un separador, aquí el guion; lpad completa el mes por la izquierda con cero hasta dos posiciones, por ejemplo cero uno para enero. Lit aporta el texto cero uno como día. To date convierte esa cadena en una fecha.}
\decir{La siguiente flecha conduce a «quarter de fecha; agrupar y sumar». Quarter devuelve el trimestre de uno a cuatro. Concat forma una etiqueta como 2020 guion T dos, uniendo el año, la letra T y el trimestre. Agrupamos por esa etiqueta y sumamos cantidades. La última flecha llega a «veinticuatro trimestres; orden cronológico»: seis años por cuatro trimestres. A diferencia del apoyo de tres c tres, aquí cada año conserva sus cuatro trimestres separados.}
\decir{La salida muestra 2019 T cuatro con siete mil ciento sesenta y cuatro; 2020 T uno con seis mil dieciocho; T dos con dos mil cuatrocientas ochenta y cinco; T tres con cuatro mil quinientas sesenta y siete. Entre esas filas, el segundo trimestre de 2020 tiene el menor acumulado. La fecha construida permite agrupar y ordenar; no representa el día real de cada desaparición ni explica la causa de la caída. Las siguientes fichas utilizarán SQL sobre vistas para otras comparaciones temporales.}
